%========================================================
% APPENDIX C - Concurrence and Reduced Density Operators
%========================================================

\chapter{Derivations for Concurrence and Reduced Density Operators}
\label{app:concurrence_reduced_density_operators}

This appendix derives the relations connecting concurrence, reduced density operators, and subsystem purity for pure bipartite two-qubit states.

\medskip

The derivation clarifies how bipartite entanglement is encoded in the mixedness of the reduced subsystems, providing the theoretical foundation for the relations used throughout Chapter~4.

%--------------------------------------------------------
% C.1 Coefficient-Matrix Representation
%--------------------------------------------------------

\section*{Coefficient-Matrix Representation}

Consider a normalized two-qubit pure state written in the computational basis as

\begin{equation}
|\psi\rangle
=
a|00\rangle
+
b|01\rangle
+
c|10\rangle
+
d|11\rangle.
\label{eq:generic_two_qubit_state}
\end{equation}

Introduce the coefficient matrix

\begin{equation}
A_\psi
=
\begin{pmatrix}
a & b\\
c & d
\end{pmatrix},
\label{eq:coefficient_matrix_definition}
\end{equation}

so that the state can be expressed compactly as

\begin{equation}
|\psi\rangle
=
\sum_{i,j=0}^{1}
(A_\psi)_{ij}
\,|i\rangle \otimes |j\rangle.
\label{eq:state_coefficient_matrix_expansion}
\end{equation}

This matrix representation provides a convenient bridge between the bipartite wavefunction and the reduced density operators.

%--------------------------------------------------------
% C.2 Reduced Density Operators
%--------------------------------------------------------

\section*{Reduced Density Operators}

The global density operator is

\begin{equation}
\rho_{AB}
=
|\psi\rangle\langle\psi|.
\label{eq:global_density_operator_appendix}
\end{equation}

Expanding Eq.~\eqref{eq:state_coefficient_matrix_expansion} yields

\begin{equation}
\rho_{AB}
=
\sum_{i,j,k,l}
(A_\psi)_{ij}
(A_\psi)^*_{kl}
\,
|i\rangle\langle k|
\otimes
|j\rangle\langle l|.
\label{eq:expanded_global_density_operator}
\end{equation}

The reduced density operator on subsystem \(A\) is defined by the partial trace

\begin{equation}
\rho_A
=
\operatorname{Tr}_B(\rho_{AB}).
\label{eq:partial_trace_appendix}
\end{equation}

Applying the partial trace to Eq.~\eqref{eq:expanded_global_density_operator} gives

\begin{equation}
\rho_A
=
\sum_{i,j,k,l}
(A_\psi)_{ij}
(A_\psi)^*_{kl}
\left(
\sum_v
\langle v|j\rangle
\langle l|v\rangle
\right)
|i\rangle\langle k|.
\label{eq:partial_trace_expansion}
\end{equation}

Using the orthogonality relation

\begin{equation}
\sum_v
\langle v|j\rangle
\langle l|v\rangle
=
\delta_{jl},
\label{eq:orthogonality_relation_appendix}
\end{equation}

one obtains

\begin{equation}
\rho_A
=
\sum_{i,k,j}
(A_\psi)_{ij}
(A_\psi)^*_{kj}
\,|i\rangle\langle k|.
\label{eq:reduced_density_operator_expanded}
\end{equation}

Therefore, in matrix form,

\begin{equation}
\rho_A
=
A_\psi A_\psi^\dagger.
\label{eq:rhoA_coefficient_matrix_relation}
\end{equation}

Similarly,

\begin{equation}
\rho_B
=
A_\psi^\dagger A_\psi.
\label{eq:rhoB_coefficient_matrix_relation}
\end{equation}

%--------------------------------------------------------
% C.3 Pure-State Concurrence
%--------------------------------------------------------

\section*{Pure-State Concurrence}

For pure two-qubit states, concurrence is defined as

\begin{equation}
C(|\psi\rangle)
=
\left|
\langle \psi |
(\sigma_y \otimes \sigma_y)
|\psi^*\rangle
\right|.
\label{eq:pure_state_concurrence_definition}
\end{equation}

Using the computational-basis representation,

\begin{equation}
|\psi\rangle
=
\begin{pmatrix}
a\\
b\\
c\\
d
\end{pmatrix},
\qquad
|\psi^*\rangle
=
\begin{pmatrix}
a^*\\
b^*\\
c^*\\
d^*
\end{pmatrix},
\label{eq:computational_basis_state_vectors}
\end{equation}

one finds

\begin{equation}
(\sigma_y \otimes \sigma_y)
|\psi^*\rangle
=
\begin{pmatrix}
- d^*\\
c^*\\
b^*\\
- a^*
\end{pmatrix}.
\label{eq:spin_flip_action}
\end{equation}

Taking the inner product yields

\begin{equation}
\langle \psi |
(\sigma_y \otimes \sigma_y)
|\psi^*\rangle
=
2(ad-bc)^*.
\label{eq:spin_flip_inner_product}
\end{equation}

Therefore,

\begin{equation}
C(|\psi\rangle)
=
2|ad-bc|.
\label{eq:pure_state_concurrence_closed_form}
\end{equation}

Using

\begin{equation}
\det(A_\psi)=ad-bc,
\label{eq:coefficient_matrix_determinant}
\end{equation}

the concurrence becomes

\begin{equation}
C(|\psi\rangle)
=
2|\det(A_\psi)|.
\label{eq:concurrence_determinant_relation}
\end{equation}

%--------------------------------------------------------
% C.4 Relation with Reduced Density Operators
%--------------------------------------------------------

\section*{Relation with Reduced Density Operators}

Using Eq.~\eqref{eq:rhoA_coefficient_matrix_relation},

\begin{equation}
\rho_A
=
A_\psi A_\psi^\dagger,
\end{equation}

one obtains

\begin{equation}
\det(\rho_A)
=
\det(A_\psi)
\det(A_\psi^\dagger)
=
|\det(A_\psi)|^2.
\label{eq:determinant_reduced_density_operator}
\end{equation}

\medskip
Substituting into Eq.~\eqref{eq:concurrence_determinant_relation} gives

\begin{equation}
C(|\psi\rangle)
=
2\sqrt{\det(\rho_A)}.
\end{equation}

This reproduces the relation introduced in Chapter~4 [see Eq.~\eqref{eq:concurrence_reduced_density_relation}].

\medskip

An identical relation holds for subsystem \( B \):

\begin{equation}
C(|\psi\rangle)
=
2\sqrt{\det(\rho_B)}.
\end{equation}

These expressions show that bipartite entanglement is directly encoded in the reduced density operators.

%--------------------------------------------------------
% C.5 Relation with Purity
%--------------------------------------------------------

\section*{Relation with Purity}

For a single-qubit density operator \( \rho \), one has

\begin{equation}
\operatorname{Tr}(\rho)=1,
\label{eq:single_qubit_trace_condition}
\end{equation}

and purity

\begin{equation}
\gamma(\rho)
=
\operatorname{Tr}(\rho^2).
\label{eq:single_qubit_purity_definition}
\end{equation}

For \(2\times2\) density operators, the determinant satisfies

\begin{equation}
\det(\rho)
=
\frac{1-\gamma(\rho)}{2}.
\label{eq:determinant_purity_relation}
\end{equation}

Substituting Eq.~\eqref{eq:determinant_purity_relation} into Eq.~\eqref{eq:concurrence_reduced_density_relation} yields

\begin{equation}
C(|\psi\rangle)
=
\sqrt{
2\left(
1-\gamma(\rho_A)
\right)
}.
\label{eq:concurrence_purity_relation}
\end{equation}

Therefore, for pure bipartite states, the mixedness of the reduced subsystem provides a direct quantitative measure of entanglement.

\medskip

Separable states produce pure reduced density operators and therefore zero concurrence, whereas maximally entangled states produce maximally mixed reduced states and maximal concurrence.